% GENERATED FILE. Edit canonical lessons, then run npm run generate:books.
% canonical-source-hash: fnv1a64:8cf793e1b785a9fe
% canonical-lessons: TA-C02-peyar, TA-C02-en, TA-C02-en-peyar, TA-C02-nii-niingal, TA-C02-enna, TA-C02-ungal-peyar-enna, TA-C02-magizhcci, TA-C02-practice

\chapter{Introducing Yourself}
\label{ch:introductions}
\hlchaptermodality{2}

\begin{chapteropening}
\textbf{By the end of this chapter:} \emph{I can introduce myself in Tamil, give my name, and understand the other person's name.}

Hold the whole introduction out loud — give your name, ask for theirs, and answer that you are pleased to meet them — using only the possession forms this chapter built.
\end{chapteropening}

\section[peyar]{\ta{பெயர்} (peyar) — \textquotedblleft{}name,\textquotedblright{} where the family is \emph{not} Indo-European}

\label{lesson:TA-C02-peyar}



\begin{quote}
\textquotedblleft{}Name\textquotedblright{} — and a word that shows Tamil at its most itself.
\end{quote}

\begin{cousinweb}
\textbf{\ta{பெயர்}} (\emph{peyar}, \textquotedblleft{}name\textquotedblright{}) is \textbf{native Dravidian}, from Proto-Dravidian \textbf{*peyar / *pēr} — and it is \emph{not} related to Sanskrit \emph{nāma} (Hindi \emph{nām}, English \emph{name}). Here the two great families part: Indo-European shares one word for \textquotedblleft{}name\textquotedblright{} from Ireland to North India, but the Dravidian south kept its own.

Across the family:

\noindent
\begin{tabularx}{\linewidth}{@{}>{\raggedright\arraybackslash}X>{\raggedright\arraybackslash}X>{\raggedright\arraybackslash}X@{}}
\toprule
\textbf{Language} & \textbf{\textquotedblleft{}Name\textquotedblright{}} & \textbf{Source} \\
\midrule
\textbf{Tamil} & \emph{peyar} & \textbf{native Dravidian} (*peyar) \\
Malayalam & \emph{pēr} & the \emph{same} Dravidian root \\
Telugu & \emph{pēru} & the \emph{same} Dravidian root \\
Kannada & \emph{hesaru} & Dravidian (*pesar, with \emph{p} $\to$ \emph{h}) \\
Hindi & \emph{nām} & Sanskrit \emph{nāman} (Indo-European) \\
English & \emph{name} & Indo-European \\
\bottomrule
\end{tabularx}

Four Dravidian tongues share \emph{peyar/pēru}; only the Indo-European languages have \emph{name/nām}.
\end{cousinweb}

\subsection*{Your turn}
Say these aloud:

\begin{itemize}
\raggedright
  \item \textquotedblleft{}peyar\textquotedblright{}
  \item the Dravidian cousins — \emph{peyar} / \emph{pēr} / \emph{pēru}
  \item what it is \emph{not} — not \emph{nām}, not \emph{name}
\end{itemize}

\begin{tcolorbox}[breakable,colback=teal!4,colframe=teal!35!black,title={Before you move on}]
Is \textbf{\ta{பெயர்}} related to English \emph{name}? (No — it's native Dravidian, \emph{*peyar}; \emph{name/nām} are Indo-European.) Which languages share \emph{peyar/pēru}? (Tamil, Malayalam, Telugu — and Kannada's \emph{hesaru}.)
\end{tcolorbox}
\section[eṉ]{\ta{என்} (eṉ) — \textquotedblleft{}my\textquotedblright{}}

\label{lesson:TA-C02-en}



\begin{quote}
\textquotedblleft{}My\textquotedblright{} — and, like \emph{peyar}, a word with no European cousin at all.
\end{quote}

\begin{cousinweb}
\textbf{\ta{என்}} (\emph{eṉ}) means \textquotedblleft{}my.\textquotedblright{} It belongs to Tamil's native Dravidian first-person family, not to the Indo-European family behind English \emph{I/my}. You will meet the related word for \textquotedblleft{}I\textquotedblright{} in its own lesson; today you only need \emph{eṉ}.
\end{cousinweb}

\subsection*{Your turn}
Say these aloud:

\begin{itemize}
\raggedright
  \item \textquotedblleft{}eṉ\textquotedblright{}
  \item \emph{eṉ} once with its meaning, \textquotedblleft{}my\textquotedblright{}
\end{itemize}

\begin{tcolorbox}[breakable,colback=teal!4,colframe=teal!35!black,title={Before you move on}]
What does \textbf{\ta{என்}} mean? (\textquotedblleft{}My.\textquotedblright{}) Does it share a root with English \emph{my}? (No — it is native Dravidian.)
\end{tcolorbox}
\section[eṉ peyar …]{\ta{என்} \ta{பெயர்} … (eṉ peyar …) — \textquotedblleft{}my name is…,\textquotedblright{} with no \textquotedblleft{}is\textquotedblright{}}

\label{lesson:TA-C02-en-peyar}



\begin{quote}
Assemble the two atoms — \emph{eṉ} and \emph{peyar} — and notice a whole word English needs that Tamil doesn't.
\end{quote}

\begin{cousinweb}
\begin{itemize}
\raggedright
  \item eṉ (my) + peyar (name) + your name.
\end{itemize}

\begin{quote}
\textbf{\ta{என்} \ta{பெயர்} …} = \textquotedblleft{}my name …\textquotedblright{} $\to$ \textbf{\textquotedblleft{}my name is…\textquotedblright{}}
\end{quote}

\emph{Eṉ peyar Arun.} $\to$ \textquotedblleft{}My name is Arun.\textquotedblright{}
\end{cousinweb}

\begin{grammarlens}[title={Grammar Lens: Tamil has no word for \textquotedblleft{}is\textquotedblright{} here}]
Look at what's missing: there is \textbf{no verb}. Tamil sets \textquotedblleft{}my name\textquotedblright{} beside \textquotedblleft{}Arun\textquotedblright{} — \emph{eṉ peyar Arun} — and the \textquotedblleft{}is\textquotedblright{} is simply understood. This is the \textbf{zero copula}: for equational sentences (X is Y), Dravidian needs no \textquotedblleft{}to be\textquotedblright{} at all — where English needs \emph{is}, Hindi needs \emph{hai}, Arabic (you'll see) also drops it. Tamil says only \textquotedblleft{}my name Arun.\textquotedblright{}
\end{grammarlens}

\subsection*{Your turn}
Say these aloud:

\begin{itemize}
\raggedright
  \item \textquotedblleft{}eṉ peyar …\textquotedblright{} with your name
  \item notice — no word for \textquotedblleft{}is\textquotedblright{}
\end{itemize}

\begin{tcolorbox}[breakable,colback=teal!4,colframe=teal!35!black,title={Before you move on}]
Give your name in Tamil. (\emph{Eṉ peyar ….}) What does Tamil leave out that English requires? (The verb \textquotedblleft{}is\textquotedblright{} — the zero copula.)
\end{tcolorbox}
\section[nī / nīṅgaḷ]{\ta{நீ} / \ta{நீங்கள்} (nī / nīṅgaḷ) — \textquotedblleft{}you,\textquotedblright{} familiar and respectful}

\label{lesson:TA-C02-nii-niingal}



\begin{quote}
To ask a name you need \textquotedblleft{}you\textquotedblright{} — and Tamil makes respect the same way French does.
\end{quote}

\begin{cousinweb}
\textbf{\ta{நீ}} (\emph{nī}, familiar \textquotedblleft{}you\textquotedblright{}) $\leftarrow$ Proto-Dravidian \textbf{*nīn} — native, unrelated to the European \textquotedblleft{}you\textquotedblright{} words (from \emph{*tū}). \textbf{\ta{நீங்கள்}} (\emph{nīṅgaḷ}) is \emph{nī} + the plural ending \textbf{-kaḷ}.
\end{cousinweb}

\begin{grammarlens}[title={Grammar Lens: respect by the plural — again}]
Tamil grades \textquotedblleft{}you\textquotedblright{} in two: \textbf{nī} (familiar — friends, children) and \textbf{nīṅgaḷ} (respectful — and literally the \textbf{plural}, \textquotedblleft{}you-all\textquotedblright{}). This is the \emph{same} politeness-by-plural as French \emph{vous}: address one respected person as \textquotedblleft{}you all.\textquotedblright{} For a first meeting, always \textbf{nīṅgaḷ}.
\end{grammarlens}

\subsection*{Your turn}
Say these aloud:

\begin{itemize}
\raggedright
  \item \textquotedblleft{}nī\textquotedblright{} (familiar), \textquotedblleft{}nīṅgaḷ\textquotedblright{} (respectful)
  \item nīṅgaḷ = nī + plural -kaḷ
\end{itemize}

\begin{tcolorbox}[breakable,colback=teal!4,colframe=teal!35!black,title={Before you move on}]
How does Tamil make \textquotedblleft{}you\textquotedblright{} respectful, and which European language uses the same trick? (By the plural — \emph{nīṅgaḷ}; French \emph{vous}.)
\end{tcolorbox}
\section[eṉṉa]{\ta{என்ன} (eṉṉa) — \textquotedblleft{}what\textquotedblright{}}

\label{lesson:TA-C02-enna}



\begin{quote}
The last piece of the question: \textquotedblleft{}what.\textquotedblright{}
\end{quote}

\begin{cousinweb}
\textbf{\ta{என்ன}} (\emph{eṉṉa}, \textquotedblleft{}what\textquotedblright{}) is native Dravidian, on the interrogative stem \textbf{*yā- / *e-} that heads Tamil's whole question family: \emph{eṉṉa} (what), \emph{yār} (who), \emph{eppo} (when), \emph{enge} (where), \emph{ēṉ} (why). Like Hindi's \emph{k-} words and English's \emph{wh-} words — but from a \emph{different}, Dravidian root.
\end{cousinweb}

\subsection*{Your turn}
Say these aloud:

\begin{itemize}
\raggedright
  \item \textquotedblleft{}eṉṉa\textquotedblright{}
  \item the Tamil question family — eṉṉa, yār, enge
\end{itemize}

\begin{tcolorbox}[breakable,colback=teal!4,colframe=teal!35!black,title={Before you move on}]
What stem heads Tamil's question-words, and is it the same as English's \emph{wh-}? (\emph{*yā-/*e-}; same \emph{idea}, different — Dravidian — root.)
\end{tcolorbox}
\section[uṅgaḷ peyar eṉṉa?]{\ta{உங்கள்} \ta{பெயர்} \ta{என்ன}? (uṅgaḷ peyar eṉṉa?) — \textquotedblleft{}what's your name?\textquotedblright{}}

\label{lesson:TA-C02-ungal-peyar-enna}



\begin{quote}
Assemble the question — your + name + what — and note, again, no \textquotedblleft{}is.\textquotedblright{}
\end{quote}

\begin{cousinweb}
\begin{itemize}
\raggedright
  \item \textbf{\ta{உங்கள்}} (\emph{uṅgaḷ}, \textquotedblleft{}your\textquotedblright{}) = the possessive of \emph{nīṅgaḷ} (\textquotedblleft{}you,\textquotedblright{} respectful).
\end{itemize}

\begin{quote}
\textbf{\ta{உங்கள்} \ta{பெயர்} \ta{என்ன}?} = literally \textquotedblleft{}\textbf{your name what?}\textquotedblright{} = \textquotedblleft{}what's your name?\textquotedblright{}
\end{quote}

Three words, \textbf{no verb} — the zero copula holds for questions too. Tamil never adds an \textquotedblleft{}is.\textquotedblright{}
\end{cousinweb}

\begin{grammarlens}[title={Grammar Lens: answer, and the familiar version}]
Answer with the sentence you built: \textbf{Eṉ peyar Arun.} The familiar \textquotedblleft{}your\textquotedblright{} swaps \emph{uṅgaḷ} for \textbf{uṉ} (from \emph{nī}): \emph{uṉ peyar eṉṉa?}
\end{grammarlens}

\subsection*{Your turn}
Say these aloud:

\begin{itemize}
\raggedright
  \item \textquotedblleft{}uṅgaḷ peyar eṉṉa?\textquotedblright{}
  \item ask, then answer — \textquotedblleft{}uṅgaḷ peyar eṉṉa?\textquotedblright{} / \textquotedblleft{}eṉ peyar …\textquotedblright{}
\end{itemize}

\begin{tcolorbox}[breakable,colback=teal!4,colframe=teal!35!black,title={Before you move on}]
What does \textbf{\ta{உங்கள்} \ta{பெயர்} \ta{என்ன}?} literally say, and what's missing versus English? (\textquotedblleft{}Your name what?\textquotedblright{}; no \textquotedblleft{}is\textquotedblright{} — zero copula.)
\end{tcolorbox}
\section[magiḻcci]{\ta{மகிழ்ச்சி} (magiḻcci) — \textquotedblleft{}joy,\textquotedblright{} and \textquotedblleft{}pleased to meet you\textquotedblright{}}

\label{lesson:TA-C02-magizhcci}



\begin{quote}
The warm close — and a sound almost unique to Tamil.
\end{quote}

\subsection*{You'll want to know: \ta{மகிழ்ச்சி}}
The full form is \emph{\textbf{uṅgaḷai sandittatil magiḻcci}} — \textquotedblleft{}\textbf{in having met you, joy}.\textquotedblright{} Say it; the extra two words are longer than you need today, and \textbf{\ta{மகிழ்ச்சி}} on its own already carries the whole meaning.

\begin{cousinweb}
\textbf{\ta{மகிழ்ச்சி}} (\emph{magiḻcci}, \textquotedblleft{}joy, gladness\textquotedblright{}) is native Dravidian, from \textbf{magiḻ} (\textquotedblleft{}to rejoice\textquotedblright{}). It carries the letter \textbf{\ta{ழ}} (\emph{ḻ}) — a retroflex glide almost unique to Tamil, with no English equivalent (the same sound in the name \emph{Tamiḻ} itself). In casual speech, a smile or the English \textquotedblleft{}nice to meet you\textquotedblright{} also does the job.
\end{cousinweb}

\subsection*{Your turn}
Say these aloud:

\begin{itemize}
\raggedright
  \item \textquotedblleft{}magiḻcci\textquotedblright{} — feel the \ta{ழ} (\emph{ḻ}) glide
  \item the fuller form — \textquotedblleft{}uṅgaḷai sandittatil magiḻcci\textquotedblright{}
\end{itemize}

\begin{tcolorbox}[breakable,colback=teal!4,colframe=teal!35!black,title={Before you move on}]
What does \textbf{\ta{மகிழ்ச்சி}} mean and from what verb? (\textquotedblleft{}Joy,\textquotedblright{} from \emph{magiḻ}, \textquotedblleft{}to rejoice.\textquotedblright{}) What rare Tamil sound does it contain? (\textbf{\ta{ழ}}, \emph{ḻ} — also in \emph{Tamiḻ}.)
\end{tcolorbox}
\section[Practice]{Chapter 2 — the introduction, whole}

\label{lesson:TA-C02-practice}



\begin{quote}
Every atom built, and every one of them native Dravidian.
\end{quote}

\subsection*{The exchange — the whole introduction}
\noindent
\begin{tabularx}{\linewidth}{@{}>{\raggedright\arraybackslash}X>{\raggedright\arraybackslash}X@{}}
\toprule
\textbf{Say this} & \textbf{English} \\
\midrule
\emph{eṉ peyar Mira.} & My name is Mira. \\
\emph{uṅgaḷ peyar eṉṉa?} & What's your name? \\
\emph{eṉ peyar Arun.} & My name is Arun. \\
\emph{magiḻcci.} & Pleased to meet you. \\
\bottomrule
\end{tabularx}

Every atom native and traced:

\begin{itemize}
\raggedright
  \item \textbf{peyar} $\leftarrow$ Dravidian \emph{*peyar} (\emph{not} English \emph{name})
  \item \textbf{eṉ} $\leftarrow$ \emph{nāṉ} (\textquotedblleft{}I\textquotedblright{})
  \item \textbf{eṉṉa} $\leftarrow$ Dravidian question-stem \emph{*yā-/*e-}
  \item \textbf{magiḻcci} $\leftarrow$ \emph{magiḻ} (\textquotedblleft{}to rejoice\textquotedblright{})
\end{itemize}

And the Tamil habit to keep: the \textbf{zero copula} — an equational sentence (\textquotedblleft{}X is Y\textquotedblright{}) needs \emph{no} word for \textquotedblleft{}is\textquotedblright{} at all. Where Hindi added \emph{hai} and English needs \emph{is}, Tamil says only \textquotedblleft{}my name Arun.\textquotedblright{}

\subsection*{Your turn}
Say these aloud:

\begin{itemize}
\raggedright
  \item the whole exchange, both voices
  \item your own introduction — \textquotedblleft{}eṉ peyar …\textquotedblright{}
  \item ask it back — \textquotedblleft{}uṅgaḷ peyar eṉṉa?\textquotedblright{}
\end{itemize}

\begin{tcolorbox}[breakable,colback=teal!4,colframe=teal!35!black,title={Before you move on}]
Give your name, ask someone else's, say you're pleased. (\emph{Eṉ peyar …. / Uṅgaḷ peyar eṉṉa? / Magiḻcci.}) What does Tamil leave out that English and Hindi require? (The verb \textquotedblleft{}is\textquotedblright{} — zero copula.)

Next chapter: \emph{eppaḍi irukkīṅgaḷ?} (\textquotedblleft{}how are you?\textquotedblright{}) and answering with \emph{nallā} — the responding cycle.
\end{tcolorbox}
